\section{解决一个实际用例：科学数据分析}

本节假设你做了某个科学实验。该科学实验产生了三张输入数据表：
每个涉及的参数 $d=2$、$d=3$、$d=4$ 各一张。这些数据表包含“自由度”
和某个精度度量“|l2_err|”。此外，它们可能还包含一些元数据
（本例中是一列“level”）。例如，$d=2$ 的数据表可能存放在
|data_d2.dat| 中，并可能包含
%
\begin{codeexample}[code only]
dof        l2_err     level
5          8.312e-02  2
17         2.547e-02  3
49         7.407e-03  4
129        2.102e-03  5
321        5.874e-04  6
769        1.623e-04  7
1793       4.442e-05  8
4097       1.207e-05  9
9217       3.261e-06  10
\end{codeexample}
%
另外两张表与之类似，我们在这里提供它们是为了便于
复现这些示例。$d=3$ 的表存放在 |data_d3.dat| 中，
内容为
%
\begin{codeexample}[code only]
dof        l2_err     level
7          8.472e-02  2
31         3.044e-02  3
111        1.022e-02  4
351        3.303e-03  5
1023       1.039e-03  6
2815       3.196e-04  7
7423       9.658e-05  8
18943      2.873e-05  9
47103      8.437e-06  10
\end{codeexample}
%
最后一张表是 |data_d4.dat|
%
\begin{codeexample}[code only]
dof        l2_err     level
9          7.881e-02  2
49         3.243e-02  3
209        1.232e-02  4
769        4.454e-03  5
2561       1.551e-03  6
7937       5.236e-04  7
23297      1.723e-04  8
65537      5.545e-05  9
178177     1.751e-05  10
\end{codeexample}

我们想做的是生成三幅图，每幅都是 dof 对 |l2_err|，放在双对数图中。
我们预期结果在双对数图中是一条直线，并且我们
关心它的斜率 $\log e(N) = -a \log(N)$，因为它刻画了我们的
实验。


\subsection{把数据读入 \TeX{}}
\label{sec:tut2:step1}

我们的第一步是把其中一张数据表读入 \PGFPlots{}。此外，
我们还想为 $x$ 轴和 $y$ 轴添加坐标轴描述，并在图的顶部加一个标题。

我们的第一版如下所示：
%
\begin{codeexample}[]
%\documentclass{article}
%\usepackage{pgfplots}
%\pgfplotsset{compat=1.5}

%\begin{document}

\begin{tikzpicture}
\begin{loglogaxis}[
    title=Convergence Plot,
    xlabel={Degrees of freedom},
    ylabel={$L_2$ Error},
]
    \addplot table {data_d2.dat};
\end{loglogaxis}
\end{tikzpicture}

%\end{document}
\end{codeexample}

我们的示例与 Section~\ref{sec:tut1:step1} 那次讲座类似：它通过
|title|、|xlabel| 和 |ylabel| 指定一些基本的坐标轴描述，并用
|\addplot table| 提供数据。唯一的区别是我们使用了
|\begin{loglogaxis}| 而不是 |\begin{axis}|，以便在两个坐标轴上
配置对数刻度。另外注意，我们省略了 |\addplot| 之后的任何
选项。如 Section~\ref{sec:tut1:step1} 所解释的，这会
让 \PGFPlots{} 查阅它的 |cycle list|，以确定一个合适的选项
列表。


\subsection{加入示例中其余的数据文件}
\label{sec:tut2:step2}

\PGFPlots{} 接受多条 |\addplot ... ;| 命令——所以我们可以直接
把其余数据文件加进来：
%
\begin{codeexample}[]
\begin{tikzpicture}
\begin{loglogaxis}[
    title=Convergence Plot,
    xlabel={Degrees of freedom},
    ylabel={$L_2$ Error},
]
    \addplot table {data_d2.dat};
    \addplot table {data_d3.dat};
    \addplot table {data_d4.dat};
\end{loglogaxis}
\end{tikzpicture}
\end{codeexample}

你也许会好奇 \PGFPlots{} 是如何为不同图选择不同线型的。你也许还会
好奇如何修改它们。其实，如果你只是写不带方括号选项的
|\addplot|，\PGFPlots{} 会为那幅特定的图自动
选择样式。这里“自动”的意思是它会查阅当前的
|cycle list|：一组预定义样式，使每条 |\addplot| 语句
得到其中一个样式。这个列表可以高度自定义。

除了 |cycle list| 之外，你也可以轻松地手动提供样式选项。如果
你写

|\addplot[|\meta{options}|] ...|，

\PGFPlots{} 将只使用 \meta{options}，并忽略它的 |cycle list|。如果
像下面这样在方括号前写一个加号

|\addplot+[|\meta{options}|] ...|，

\PGFPlots{} 会把 \meta{options} 追加到自动分配的循环
列表之后。


\subsection{加上图例和网格}
\label{sec:tut2:step3}

图例是一段文字标签，说明各幅图是什么。可以使用 |legend entries|
键为一条或多条 |\addplot| 语句提供图例：
%
\begin{codeexample}[]
\begin{tikzpicture}
\begin{loglogaxis}[
    title=Convergence Plot,
    xlabel={Degrees of freedom},
    ylabel={$L_2$ Error},
    grid=major,
    legend entries={$d=2$,$d=3$,$d=4$},
]
    \addplot table {data_d2.dat};
    \addplot table {data_d3.dat};
    \addplot table {data_d4.dat};
\end{loglogaxis}
\end{tikzpicture}
\end{codeexample}
%
这里，我们指定了一个逗号分隔的文字标签列表，每条 |\addplot|
指令一个。注意文字标签中数学模式的用法。另外注意，
如果你的某个标签包含逗号，就必须用花括号把该项
括起来。例如，我们本可以使用
|legend entries={{$d=2$},{$d=3$},{$d=4$}}|——\PGFPlots{} 用这些花括号
界定参数，之后再去掉它们（顺便说一句，这对任何选项都
成立）。

我们的示例还包含网格线，为此使用了 |grid=major| 键。它
在所有坐标轴中激活主网格线。

你也许会好奇这些文字标签是如何对应到 |\addplot| 指令的。其实，
它们按序号对应。第一个标签赋给第一幅图，第二个
标签赋给第二幅图，依此类推。如果你给某幅图加上
|forget plot| 选项（例如使用 |\addplot+[forget plot]|），就可以把该图
排除在这个计数之外。这样的图同时被排除在循环列表和
图例之外。


\subsection{加一条选定的拟合线}
\label{sec:tut2:step4}

有时，需要对输入样本计算线性回归直线。假设我们想为
第四张数据表（|data_d4.dat|）中的数据计算一条拟合线。
不过，我们假设图中有趣的部分出现在自由度数目达到某个
渐近极限时（即非常大时）。因此，在计算拟合线时，我们想
给前几个点赋予较高的
不确定性。

\PGFPlots{} 支持把表输入与数学表达式结合起来（注意，
你也可以直接写纯数学表达式，不过这超出了本例的
范围）。本例中，我们利用这个特性创建一个
全新的列——线性回归直线：
%
\begin{codeexample}[]
%\usepackage{pgfplotstable}
%...
\begin{tikzpicture}
\begin{loglogaxis}[
    title=Convergence Plot,
    xlabel={Degrees of freedom},
    ylabel={$L_2$ Error},
    grid=major,
    legend entries={$d=2$,$d=3$,$d=4$},
]
    \addplot table {data_d2.dat};
    \addplot table {data_d3.dat};
    \addplot table {data_d4.dat};
    \addplot table [
      x=dof,
      y={create col/linear regression={y=l2_err,
       variance list={1000,800,600,500,400,200,100}}}]
        {data_d4.dat};
\end{loglogaxis}
\end{tikzpicture}
\end{codeexample}

注意，我们又加了一个宏包：|pgfplotstable|。它允许对
表做后处理（等等。它还有一个功能强大的表排版工具箱，
能根据你输入的 CSV 文件对数字进行四舍五入和格式化）。

这里，我们给坐标轴加了第四幅图。第一幅图同样是一条
|\addplot table| 语句——我们看到它像前一幅图一样加载数据文件
|data_d4.dat|。不过，它带有一些控制坐标输入的特殊键：
|x=dof| 表示从名为“dof”的列加载 x 坐标。这与我们其他所有图
本质上相同（因为“dof”列就是第一列）。它还使用
|y={create col/...}|。这条长长的语句定义了一个全新的列。
|create col/linear regression| 前缀是一个键，凡是
可以生成新表列的地方都能使用它。一旦表被
第一次查询，该语句就会被求值，然后应用于所有后续行。
|create col/linear regression| 的参数表包含用于回归直线的
函数值列名 |y=l2_err|（x 参数正如你猜对的那样，由 |x=dof|
推出）。|variance list| 选项是可选的。我们用它给
前几个输入点赋予方差
（不确定性）。更确切地说：第一个遇到的数据点得到方差
1000，第二个 800，第三个 600，依此类推。方差的个数不必与
点数匹配；\PGFPlots{} 只是把它们与最先遇到的
坐标配对。

注意，由于我们的 |legend entries| 键只含三个值，
回归直线没有图例项。如果想加，我们可以轻松加上一个。我们
也可以使用 |\addplot+[forget plot] table[...]| 来如上所述显式
抑制图例的生成。

每当 \PGFPlots{} 遇到数学表达式时，它就使用内置的
浮点单元。因此，它的数据范围非常大——精度也
相当不错。


\subsection{用 \Tikz{} 添加标注：斜率三角形}
\label{sec:tut2:step5}

数据往往需要解读——你可能想在图中突出特定的
项。这种“突出特定项”需要画进
坐标轴内部，而且需要高度的灵活性。\Tikz{} 的用户会说
\Tikz{} 是理所当然的选择——确实如此。

在我们的用例中，我们关心斜率。我们可能想比较
不同实验的斜率。我们还可能想展示选定的斜率
绝对值。

这里，我们用 \Tikz{} 往 \PGFPlots{} 坐标轴中添加自定义标注。我们
选择了一种特定类型的自定义标注：我们想在
线图上标记两个点。一种做法是确定确切坐标，
并在该坐标处放置一个图形元素（这可以用
|\draw ... (1e4,1e-5) ... ;| 实现）。另一种（可能更简单的）做法是使用
|pos| 功能来标识“线开始后 25\%”的位置。

基于 Section~\ref{sec:tut2:step4} 的结果，我们得到
%
\begin{codeexample}[]
%\usepackage{pgfplotstable}
%...
\begin{tikzpicture}
\begin{loglogaxis}[
    title=Convergence Plot,
    xlabel={Degrees of freedom},
    ylabel={$L_2$ Error},
    grid=major,
    legend entries={$d=2$,$d=3$,$d=4$},
]
    \addplot table {data_d2.dat};
    \addplot table {data_d3.dat};
    \addplot table {data_d4.dat};
    \addplot table [
      x=dof,
      y={create col/linear regression={y=l2_err,
       variance list={1000,800,600,500,400,200,100}}}]
        {data_d4.dat}
            % save two points on the regression line
            % for drawing the slope triangle
            coordinate [pos=0.25] (A)
            coordinate [pos=0.4]  (B)
   ;
    % save the slope parameter:
    \xdef\slope{\pgfplotstableregressiona}

    % draw the opposite and adjacent sides
    % of the triangle
    \draw (A) -| (B)
         node [pos=0.75,anchor=west]
         {\pgfmathprintnumber{\slope}};
\end{loglogaxis}
\end{tikzpicture}
\end{codeexample}

这个示例已经相当复杂，因为我们每一步都在增加复杂度。
在深入细节之前，先看一个更简单的例子：
%
\begin{codeexample}[]
\begin{tikzpicture}
\begin{loglogaxis}
    \addplot table {data_d2.dat}
       coordinate [pos=0.25] (A)
       coordinate [pos=0.4]  (B)
    ;

    \draw [-stealth] (A) -| (B);

    \node [pin=0:Special.] at (769,1.623e-04) {};
\end{loglogaxis}
\end{tikzpicture}
\end{codeexample}
%
这里，我们看到 \PGFPlots{} 提供的两种标注概念：第一种是
在一条 |\addplot| 命令之后（但在
结束分号之前）直接插入绘图命令。第二种是添加标准的 \Tikz{}
命令，但使用指定的 \PGFPlots{} 坐标。两者都是
\Tikz{} 的概念。第一种
正是我们这里想要的：我们想标识出线上
“某处”的两个坐标。本例中，我们定义两个命名坐标：坐标 $A$ 位于
线的 25\% 处，坐标 $B$ 位于线的 40\% 处。然后，我们用
\verb#\draw (A) -| (B)# 在这两点之间画一个三角形。第二种
只在你事先知道某些绝对坐标时才有用。

回到我们最初带回归直线的做法，可以看到它
使用的是第一种概念：它在 |\addplot| 之后、
结束分号之前引入命名坐标。语句 |\xdef\slope| 引入一个新的
宏。它包含 |\pgfplotstableregressiona| 的值（因“eXpanded DEFinition”而被展开），
即回归直线的斜率。除了斜率三角形之外，我们还加了一个节点，
在其中用 |\pgfmathprintnumber| 排版该值。

注意，上面的例子其实是一个“顺利的情形”：很容易发生
坐标轴环境内部添加的标签被裁剪掉的
情况：
%
\begin{codeexample}[]
\begin{tikzpicture}
\begin{loglogaxis}[
    tiny,
]
    \addplot table {data_d2.dat}
        node [pos=1,pin=0:Special.] {}
    ;
\end{loglogaxis}
\end{tikzpicture}
\end{codeexample}
%
上面的例子把 |pos| 标签放置与节点的标签结合了起来。
注意，小号样式 |tiny| 安装了一个 \PGFPlots{} 预设，它更适合
非常小的图——它是众多受支持的缩放
参数之一。这里的问题显而易见：我们额外节点的文字被裁剪
掉了。根据你的数据，这里有几种
解决方案：
%
\begin{itemize}
    \item 使用 |clip=false| 完全禁用绘图路径的裁剪，
    \item 使用 |clip mode=individual| 只对绘图路径启用裁剪，
    \item 把节点画在坐标轴环境之外、但仍在
        picture 环境内部。
\end{itemize}
%
第一种做法对大多数图都相当有效：
%
\begin{codeexample}[]
\begin{tikzpicture}
\begin{loglogaxis}[
    tiny,
    clip=false,
]
    \addplot table {data_d2.dat}
        node [pos=1,pin=0:Special.] {}
    ;
\end{loglogaxis}
\end{tikzpicture}
\end{codeexample}
%
注意，这种把节点放在结束分号之前的做法意味着节点会
继承坐标轴线的样式和颜色。


\subsection{小结}

我们学习了如何定义一个（对数）坐标轴，以及如何指定基本的
坐标轴描述。我们还再次看到了如何使用一条或多条
|\addplot table| 命令把表数据读入 \PGFPlots{}。我们简要了解了
回归分析和 \Tikz{} 绘图标注。

我们还遇到了 |tiny| 样式，它是自定义坐标轴
尺寸的方式之一。其他方式还有 |width|、|height|，以及其他预定义尺寸
样式如 |normalsize|、|small| 或 |footnotesize|，还有两种不同的缩放
模式 |/pgfplots/scale| 和 |/tikz/scale|（前者只缩放坐标轴，
后者同时缩放标签）。

下一步可能包括：如何用线图可视化函数，如何正确对齐
相邻的图形（即使坐标轴描述各不相同），如何使用
\PGFPlots{} 的散点图，或如何绘制二元函数。

%The steps are explained in all detail in the reference
%Chapter~\ref{cha:pgfplots:reference}.
